\setcounter{section}{11}

  \zwischenueberschrift{Übungsaufgaben}

\inputaufgabe
{}
{

Zeige, dass das
\definitionsverweis {Produkt}{}{}

\mathl{M \times N}{} von zwei
\definitionsverweis {differenzierbaren Mannigfaltigkeiten}{}{}
\mathkor {} {M} {und} {N} {}
selbst wieder eine differenzierbare Mannigfaltigkeit ist.

}
{} {}

\inputaufgabe
{}
{

Es seien
\mathkor {} {M_1 \subseteq N_1} {und} {M_2 \subseteq N_2} {}
\definitionsverweis {abgeschlossene Untermannigfaltigkeiten}{}{.} Zeige, dass ihr
\definitionsverweis {Produkt}{}{}
\mathl{M_1 \times M_2}{} eine abgeschlossene Untermannigfaltigkeit von
\mathl{N_1 \times N_2}{} ist.

}
{} {}

\inputaufgabe
{}
{

Beschreibe die Karten auf dem
\definitionsverweis {Torus}{}{}

\mathl{S^1 \times S^1}{,} die von den stereographischen Projektionen herrühren.

}
{} {}

\inputaufgabe
{}
{

Zeige, dass
\mathl{\R^2 \setminus \{0\}}{}
\definitionsverweis {diffeomorph}{}{}
zu einem Produkt aus eindimensionalen
\definitionsverweis {Mannigfaltigkeiten}{}{}
ist.

}
{} {}

\inputaufgabe
{}
{

Zeige, dass das
\definitionsverweis {Produkt}{}{}

\mathl{M \times N}{} von zwei
\definitionsverweis {wegzusammenhängenden}{}{} \definitionsverweis {differenzierbaren Mannigfaltigkeiten}{}{}
\mathkor {} {M} {und} {N} {}
wieder wegzusammenhängend ist.

}
{} {}

\inputaufgabe
{}
{

Es sei $M$ eine
\definitionsverweis {differenzierbare Mannigfaltigkeit}{}{}
und
\maabbeledisp {\varphi} { M } { M \times M
} { x } {(x,x)
} {,}
die
\definitionsverweis {Diagonalabbildung}{}{}
in das
\definitionsverweis {Produkt}{}{}

\mathl{M \times M}{.} Zeige, dass die
\definitionsverweis {Diagonale}{}{}
$\varphi(M)$ eine
\definitionsverweis {abgeschlossene Untermannigfaltigkeit}{}{}
ist.

}
{} {}

\inputaufgabe
{}
{

Es sei $M$ eine
\definitionsverweis {differenzierbare Mannigfaltigkeit}{}{.}
Zeige, dass die Vertauschungsabbildung
\maabbeledisp {} {M \times M} { M \times M
} {(P,Q)} {(Q,P)
} {,}
ein
\definitionsverweis {Diffeomorphismus}{}{}
ist.

}
{} {}

\inputaufgabe
{}
{

Beschreibe den Torus
\mathl{S^1 \times S^1}{} als
\definitionsverweis {Rotationsmenge}{}{}
im $\R^3$.

}
{} {}

\inputaufgabe
{}
{

Es sei

\mavergleichskette
{\vergleichskette
{ R
}
{ > }{ 0
}
{  }{
}
{    }{
}
{   }{
}
}
{}{}{}
und betrachte die Abbildung
\maabbeledisp {} { \R^3} { \R
} { (x,y,z) } { \left(  \sqrt{x^2+y^2}- R  \right)^2 +z^2
} {.}
Bestimme die
\definitionsverweis {regulären Punkte}{}{}
der Abbildung und die Gestalt der Faser über

\mavergleichskette
{\vergleichskette
{ s
}
{ \in }{ \R
}
{  }{
}
{    }{
}
{   }{
}
}
{}{}{.}
Wie ändert sich die Gestalt beim Übergang von

\mavergleichskette
{\vergleichskette
{ \sqrt{s}
}
{ < }{ R
}
{  }{
}
{    }{
}
{   }{
}
}
{}{}{}
zu

\mavergleichskette
{\vergleichskette
{ \sqrt{s}
}
{ > }{ R
}
{  }{
}
{    }{
}
{   }{
}
}
{}{}{?}

}
{} {}

\inputaufgabegibtloesung
{}
{

Erstelle eine Animation, die
[[Torus/Faserrealisierung/Andere Fasern/Aufgabe|Aufgabe 11.9]]
illustriert.

}
{} {}

\inputaufgabe
{}
{

Definiere die Abbildung
\maabbdisp {} {S^1 \times S^1} { S^2
} {,}
die zu einem Winkelpaar
\mathl{(\alpha, \beta)}{} die erste Komponente als Äquatorpunkt interpretiert und von dort aus mit der zweiten Komponente auf dem Großkreis Richtung Norden wandert. Ist die Abbildung
\definitionsverweis {differenzierbar}{}{?}
Wie sehen die Fasern der Abbildung aus?

}
{} {}

\inputaufgabe
{}
{

Man gebe ein heuristisches Argument, dass die Einheitssphäre
\mathl{S^2}{} und der Torus
\mathl{S^1 \times S^1}{} nicht
\definitionsverweis {homöomorph}{}{}
sind.

}
{} {}

\inputaufgabe
{}
{

Zu welcher
\definitionsverweis {differenzierbaren Mannigfaltigkeit}{}{}
ist
\mathl{S^1 \times S^1 \setminus \triangle}{,} also der Torus ohne die
\definitionsverweis {Diagonale}{}{,}
\definitionsverweis {diffeomorph}{}{?}

}
{} {}

\inputaufgabegibtloesung
{}
{

Es sei $X$ ein
\definitionsverweis {Torus}{}{.}
Man gebe eine surjektive
\definitionsverweis {differenzierbare Abbildung}{}{}
\maabbdisp {\varphi} {\R^2} {X
} {}
derart an, dass auch die
\definitionsverweis {Tangentialabbildung}{}{}
\maabbdisp {T_P(\varphi)} { T_P\R^2 } { T_{\varphi(P)}X
} {}
in jedem Punkt

\mavergleichskette
{\vergleichskette
{ P
}
{ \in }{ \R^2
}
{  }{
}
{    }{
}
{   }{
}
}
{}{}{}
surjektiv ist.

}
{} {}

In den folgenden Aufgaben wird das relative Produkt von Mengen und topologischen Räumen eingeführt.

\inputaufgabe
{}
{

Es seien
\mathl{L_1,L_2,M}{} Mengen und
\maabbdisp {p_1} {L_1 } { M
} {}
und
\maabbdisp {p_2} {L_2 } { M
} {}
Abbildungen. Wir definieren

\mavergleichskettedisp
{\vergleichskette
{L_1 \times_M L_2
}
{ \defeq} { \left\{ (x_1,x_2)  \mid   p_1(x_1) = p_2(x_2)         \right\}
}
{ \subseteq} { L_1 \times L_2
}
{ } {
}
{ } {
}
}
{}{}{.}
\aufzaehlungzwei {Zeige, dass es ein kommutatives Diagramm

\mathdisp {\begin{matrix} L_1 \times_M L_2  & \stackrel{  }{\longrightarrow} &  L_1  &  \\ \downarrow & & \downarrow  & \\  L_2  & \stackrel{  }{\longrightarrow} &  M & \! \! \! \!\!\!\!\!   \\ \end{matrix}} {  }

gibt.
} {Es sei $T$ eine weitere Menge und
\maabb {\psi_1} { T } { L_1
} {}
und
\maabb {\psi_2} { T } { L_2
} {}
Abbildungen mit

\mavergleichskettedisp
{\vergleichskette
{ p_1 \circ \psi_1
}
{ =} { p_2 \circ \psi_2
}
{ } {
}
{ } {
}
{ } {}
}
{}{}{.}
Zeige, dass es eine eindeutig bestimmte Abbildung
\maabbdisp {\psi} { T } { L_1 \times_M L_2
} {}
derart gibt, dass die Projektionen auf
\mathkor {} {L_1} {bzw.} {L_2} {}
mit $\psi_1,\psi_2$ übereinstimmen.
}

}
{} {}

\inputaufgabe
{}
{

Es seien
\mathl{L_1,L_2,M}{}
\definitionsverweis {topologische Räume}{}{}
und
\maabbdisp {p_1} {L_1 } { M
} {}
und
\maabbdisp {p_2} {L_2 } { M
} {}
\definitionsverweis {stetige Abbildungen}{}{.}
Wir definieren

\mavergleichskettedisp
{\vergleichskette
{L_1 \times_M L_2
}
{ \defeq} { \left\{ (x_1,x_2)  \mid   p_1(x_1) = p_2(x_2)         \right\}
}
{ \subseteq} { L_1 \times L_2
}
{ } {
}
{ } {
}
}
{}{}{}
mit der
\definitionsverweis {induzierten Topologie}{}{.}
\aufzaehlungzwei {Zeige, dass es ein kommutatives Diagramm

\mathdisp {\begin{matrix} L_1 \times_M L_2  & \stackrel{  }{\longrightarrow} &  L_1  &  \\ \downarrow & & \downarrow  & \\  L_2  & \stackrel{  }{\longrightarrow} &  M & \! \! \! \!\!\!\!\!   \\ \end{matrix}} {  }

mit stetigen Abbildungen gibt.
} {Es sei $T$ ein weiterer topologischer Raum und
\maabb {\psi_1} { T } { L_1
} {}
und
\maabb {\psi_2} { T } { L_2
} {}
stetige Abbildungen mit

\mavergleichskettedisp
{\vergleichskette
{ p_1 \circ \psi_1
}
{ =} { p_2 \circ \psi_2
}
{ } {
}
{ } {
}
{ } {}
}
{}{}{.}
Zeige, dass es eine eindeutig bestimmte stetige Abbildung
\maabbdisp {\psi} { T } { L_1 \times_M L_2
} {}
derart gibt, dass die Projektionen auf
\mathkor {} {L_1} {bzw.} {L_2} {}
mit $\psi_1,\psi_2$ übereinstimmen.
}

}
{} {}

In der Situation der beiden vorstehenen Aufgaben bezeichnet man $L_1 \times_M L_2$
\zusatzklammer {mit der Projektion auf $L_2$} {} {}
auch als $p_2^*L_1$ bzw. als
\mathl{p_1^*L_2}{,} die Vorstellung ist dabei, dass man das Objekt
\maabb {p_1} {L_1} {M
} {}
längs $p_2$ auf die neue Basis $L_2$ zurückzieht.

\inputaufgabe
{}
{

Es sei
\maabb {p} { Y } { X
} {}
eine
\definitionsverweis {stetige Abbildung}{}{}
zwischen den
\definitionsverweis {topologischen Räumen}{}{}
\mathkor {} {X} {und} {Y} {.}
Zeige, dass für das
\definitionsverweis {relative Produkt}{}{}
die Beziehung

\mavergleichskettedisp
{\vergleichskette
{ X \times_X Y
}
{ \cong} { Y
}
{ } {
}
{ } {
}
{ } {
}
}
{}{}{}
gilt
\zusatzklammer {wobei
\maabb {} { X } { X
} {}
die Identität sei} {} {.}

}
{} {}

\inputaufgabe
{}
{

Es sei
\maabb {p} { Y } { X
} {}
eine
\definitionsverweis {stetige Abbildung}{}{}
zwischen den
\definitionsverweis {topologischen Räumen}{}{}
\mathkor {} {X} {und} {Y} {.}
Es sei

\mavergleichskette
{\vergleichskette
{ x
}
{ \in }{ X
}
{  }{
}
{    }{
}
{   }{
}
}
{}{}{}
ein Punkt mit der zugehörigen Abbildung
\maabbdisp {\iota} {\{x\}} { X
} {.}
Zeige, dass für das
\definitionsverweis {relative Produkt}{}{}
die Beziehung

\mavergleichskettedisp
{\vergleichskette
{ \{x\}  \times_X Y
}
{ \cong} { Y_x
}
{ } {
}
{ } {
}
{ } {
}
}
{}{}{}
gilt, also das relative Produkt mit der
\definitionsverweis {Faser}{}{}
übereinstimmt.

}
{} {}

\inputaufgabe
{}
{

Es seien
\mathkor {} {X,U} {und} {V} {}
\definitionsverweis {topologische Räume}{}{}
und seien
\maabb {p_1} {X \times U} {X
} {}
und
\maabb {p_2} {X \times V} {X
} {}
die kanonischen Projektionen nach $X$. Zeige, dass für das
\definitionsverweis {relative Produkt}{}{}
die Beziehung

\mavergleichskettedisp
{\vergleichskette
{ \left(  X \times U  \right) \times_X \left(  X \times V  \right)
}
{ \cong} { X \times U \times V
}
{ } {
}
{ } {
}
{ } {
}
}
{}{}{}
gilt.

}
{} {}

\inputaufgabe
{}
{

Es seien
\maabb {f,g} {\R} {\R
} {}
\definitionsverweis {stetige Funktionen}{}{.}
\aufzaehlungdrei{Zeige, dass für das
\definitionsverweis {relative Produkt}{}{}
\zusatzklammer {die Funktionen treten in der Notation nicht explizit auf} {} {}

\mavergleichskettedisp
{\vergleichskette
{ \R \times_\R \R
}
{ =} { \left\{ (x,y) \in \R^2   \mid   f(x) = g(y)         \right\}
}
{ } {
}
{ } {
}
{ } {
}
}
{}{}{}
gilt.
}{Skizziere
\mathl{\R \times_\R \R}{} für eine Auswahl an Funktionen
\mathl{(f,g)}{.}
}{Bestimme in (2), ob
\mathl{\R \times_\R \R}{}
\definitionsverweis {singuläre Punkte}{}{}
besitzt.
}

}
{} {}

\inputaufgabe
{}
{

Es seien
\mathkor {} {X,Y} {und} {Z} {}
\definitionsverweis {topologische Räume}{}{}
und seien
\maabb {\varphi} { Y } { X
} {}
und
\maabb {p} { Z } { X
} {}
\definitionsverweis {stetige Abbildungen}{}{.}
Es sei
\maabbdisp {p_Y} {Y \times_XZ } { Y
} {.}
Zeige, dass ein
\definitionsverweis {stetiger Schnitt}{}{}
\maabb {s} { Y } { Y \times_XZ
} {}
das gleiche ist wie eine stetige Abbildung
\maabb {t} { Y } { Z
} {}
mit

\mavergleichskette
{\vergleichskette
{ p \circ t
}
{ = }{ \varphi
}
{  }{
}
{    }{
}
{   }{
}
}
{}{}{.}

}
{} {}

\inputaufgabe
{}
{

Es seien
\mathl{X,Y,Z}{} und $X'$
\definitionsverweis {topologische Räume}{}{}
und es liege ein kommutatives Diagramm von
\definitionsverweis {stetigen Abbildungen}{}{}

\mathdisp {\begin{matrix}  & Y& \stackrel{}{\longrightarrow} &Z  \\  &  & \searrow& \downarrow   \\  & X' & \stackrel{\varphi}{  \longrightarrow } & X    \end{matrix}} {  }

vor. Es sei

\mavergleichskette
{\vergleichskette
{ Y'
}
{ = }{ \varphi^*Y
}
{ = }{ Y \times_X X'
}
{    }{
}
{   }{
}
}
{}{}{}
und

\mavergleichskette
{\vergleichskette
{ Z'
}
{ = }{ \varphi^*Z
}
{ = }{ Z \times_X X'
}
{    }{
}
{   }{
}
}
{}{}{.}
Zeige, dass ein kommutatives Diagramm

\mathdisp {\begin{matrix}  & Y' & \stackrel{}{\longrightarrow} & Z'  \\  &  & \searrow& \downarrow   \\  &  &  & X'    \end{matrix}} {  }

vorliegt.

}
{} {}

\inputaufgabe
{}
{

Es seien
\mathkor {} {X,Y,Z} {und} {W} {}
\definitionsverweis {topologische Räume}{}{}
und seien
\maabb {} { W } { X
} {}
und

\mathdisp {Z \stackrel{\varphi}{\longrightarrow} Y\stackrel{\psi}{\longrightarrow} X} {  }

\definitionsverweis {stetige Abbildungen}{}{.}
Zeige, dass das
\definitionsverweis {relative Produkt}{}{}
\zusatzklammer {beziehungsweise der Rückzug} {} {}
die Beziehung

\mavergleichskettedisp
{\vergleichskette
{ \varphi^* \left(  \psi^*W  \right)
}
{ =} { Z \times_Y \left(  Y \times_X W  \right)
}
{ =} { Z \times_XW
}
{ =} { ( \psi \circ \varphi)^*W
}
{ } {
}
}
{}{}{}
erfüllt.

}
{} {}

\inputaufgabe
{}
{

Betrachte die
\definitionsverweis {Kreislinie}{}{}
$S^1$. Definiere eine \stichwort {differenzierbare Gruppenstruktur} {} auf $S^1$, also ein neutrales Element
\mathl{P \in S^1}{,} eine
\definitionsverweis {differenzierbare Abbildung}{}{}
\maabbeledisp {\alpha} {S^1} {S^1
} { x } { \alpha (x)
} {,}
und eine differenzierbare Abbildung
\maabbeledisp {} {T = S^1 \times S^1} {S^1
} {(x,y)} { \varphi(x,y)
} {,}
derart, dass $S^1$ mit diesen Daten zu einer
\definitionsverweis {kommutativen Gruppe}{}{}
wird.

}
{} {}

\inputaufgabe
{}
{

Betrachte die
\definitionsverweis {allgemeine lineare Gruppe}{}{}

\mavergleichskette
{\vergleichskette
{G
}
{ = }{ \operatorname{GL}_{  n   } \! \left(  \R  \right)
}
{ \subseteq }{ \R^{n^2}
}
{    }{
}
{   }{
}
}
{}{}{}
als
\definitionsverweis {offene Untermannigfaltigkeit}{}{}
des $\R^{n^2}$. Definiere eine \stichwort {differenzierbare Gruppenstruktur} {} auf $G$, also ein neutrales Element
\mathl{P \in G}{,} eine
\definitionsverweis {differenzierbare Abbildung}{}{}
\maabbeledisp {\alpha} { G } { G
} { x } { \alpha (x)
} {,}
und eine differenzierbare Abbildung
\maabbeledisp {\varphi} {G \times G} {G
} {(x,y)} { \varphi(x,y)
} {,}
derart, dass $G$ mit diesen Daten zu einer
\definitionsverweis {Gruppe}{}{}
wird.

}
{} {}

\inputaufgabe
{}
{

Zeige, dass die
\definitionsverweis {Abbildung}{}{}
\maabbeledisp {\varphi} {S^1 \times \R } { S^1 \times \R
} {(x_1 , x_2;t)} {(-x_1 , -x_2; -t)
} {,}
ein
\definitionsverweis {fixpunktfreier}{}{} \definitionsverweis {Diffeomorphismus}{}{}
ist, der zu sich selbst invers ist.

}
{} {}

\inputaufgabe
{}
{

Es sei

\mavergleichskette
{\vergleichskette
{ D
}
{ = }{ x^2 \partial_x +  \left(  x-y^2  \right) \partial_y
}
{  }{
}
{    }{
}
{   }{
}
}
{}{}{.}
Berechne

\mathdisp {D  \left(  4x^2-xy+5y^3  \right)} { . }

}
{} {}

\inputaufgabe
{}
{

Es sei

\mavergleichskette
{\vergleichskette
{ D
}
{ = }{ \sum_{j = 1}^n g_j \partial_j
}
{  }{
}
{    }{
}
{   }{
}
}
{}{}{}
ein Differentialoperator auf einer offenen Menge

\mavergleichskette
{\vergleichskette
{ U
}
{ \subseteq }{ \R^n
}
{  }{
}
{    }{
}
{   }{
}
}
{}{}{.}
Die Auswertung von $D$ an welchen differenzierbaren Funktionen gibt die Koeffizientenfunktionen $g_j$  aus?

}
{} {}

\inputaufgabe
{}
{

Es sei

\mavergleichskette
{\vergleichskette
{ D
}
{ = }{ \sum_{j = 1}^n g_j \partial_j
}
{  }{
}
{    }{
}
{   }{
}
}
{}{}{}
ein Differentialoperator auf einer offenen Menge

\mavergleichskette
{\vergleichskette
{ U
}
{ \subseteq }{ \R^n
}
{  }{
}
{    }{
}
{   }{
}
}
{}{}{.}
Zeige, dass $D$ ein Differentialoperator im Sinne der
[[Mannigfaltigkeit/Differentialoperator/Erste Ordnung/Definition|Definition 11.4]]
ist.

}
{} {}

\inputaufgabe
{}
{

Es sei $M$ eine
\definitionsverweis {differenzierbare Mannigfaltigkeit}{}{}
und $D$ ein differenzierbares
\definitionsverweis {Vektorfeld}{}{}
auf $M$. Zeige

\mavergleichskette
{\vergleichskette
{ [ D,D]
}
{ = }{ 0
}
{  }{
}
{    }{
}
{   }{
}
}
{}{}{.}

}
{} {}

\inputaufgabe
{}
{

Es sei $M$ eine
$C^3$-\definitionsverweis {differenzierbare Mannigfaltigkeit}{}{.}
Zeige, dass für
$C^2$-\definitionsverweis {Vektorfelder}{}{}
$D,E,F$ auf $M$ die sogenannte \stichwort {Jacobi-Identität} {}

\mavergleichskettedisp
{\vergleichskette
{ [ [D,E],F] + [ [E,F],D] + [ [F,D],E]
}
{ =} { 0
}
{ } {
}
{ } {
}
{ } {
}
}
{}{}{}
gilt.

}
{} {}

Es sei $V$ ein
\definitionsverweis {Vektorraum}{}{}
über einem
\definitionsverweis {Körper}{}{}
$K$ und sei
\maabbeledisp {} {V \times V} { V
} {(f,g)} { [f,g]
} {,}
eine
\definitionsverweis {Verknüpfung}{}{}
auf $V$. Man nennt
\mathl{(V, [-,-])}{} eine
\definitionswort {Lie-Algebra}{,}
wenn die folgenden drei Bedingungen erfüllt sind.
\aufzaehlungdrei{ $[-,-]$ ist
\definitionsverweis {bilinear}{}{.}
}{ $[-,-]$ erfüllt die sogenannte \stichwort {Jacobi-Identität} {,} d.h. für

\mavergleichskette
{\vergleichskette
{ u,v,w
}
{ \in }{ V
}
{  }{
}
{    }{
}
{   }{
}
}
{}{}{}
gilt

\mavergleichskettedisp
{\vergleichskette
{ [ [u,v],w] + [ [v,w],u] + [ [w,u],v]
}
{ =} { 0
}
{ } {
}
{ } {
}
{ } {
}
}
{}{}{.}
}{Es ist

\mavergleichskette
{\vergleichskette
{ [v,v]
}
{ = }{ 0
}
{  }{
}
{    }{
}
{   }{
}
}
{}{}{}
für alle

\mavergleichskette
{\vergleichskette
{ v
}
{ \in }{ V
}
{  }{
}
{    }{
}
{   }{
}
}
{}{}{.}
}

\inputaufgabe
{}
{

Es sei $M$ eine
$C^\infty$-\definitionsverweis {Mannigfaltigkeit}{}{}
und sei $V$ der Raum aller
$C^\infty$-\definitionsverweis {Vektorfelder}{}{}
auf $M$. Zeige, dass  $V$ mit der
\definitionsverweis {Lie-Klammer}{}{}
von Vektorfeldern eine
\definitionsverweis {Lie-Algebra}{}{}
ist.

}
{} {}

\inputaufgabe
{}
{

Es sei $V$ ein
$K$-\definitionsverweis {Vektorraum}{}{}
über einem
\definitionsverweis {Körper}{}{}
$K$ und sei

\mavergleichskettedisp
{\vergleichskette
{ L
}
{ =} { \operatorname{End}_{  } \, ( V )
}
{ } {
}
{ } {
}
{ } {
}
}
{}{}{}
die Menge aller linearen
$K$-\definitionsverweis {Endomorphismen}{}{}
von $V$ nach $V$. Zeige, dass $L$, versehen mit

\mavergleichskettedisp
{\vergleichskette
{ [f,g]
}
{ \defeq} { f \circ g -g \circ f
}
{ } {
}
{ } {
}
{ } {
}
}
{}{}{,}
eine
\definitionsverweis {Lie-Algebra}{}{}
ist.

}
{} {}

  \zwischenueberschrift{Aufgaben zum Abgeben}

\inputaufgabe
{5}
{

Es sei

\mavergleichskette
{\vergleichskette
{ 0
}
{ < }{ r
}
{ < }{ R
}
{    }{
}
{   }{
}
}
{}{}{}
und sei

\mavergleichskettedisp
{\vergleichskette
{ T
}
{ =} { \left\{ (x,y,z) \in \R^3  \mid   \left(  \sqrt{x^2+y^2} -R  \right)^2 +z^2 = r^2         \right\}
}
{ } {
}
{ } {
}
{ } {
}
}
{}{}{.}
Zeige, dass die Abbildung
\maabbeledisp {} { S^1 \times S^1 } { T
} { ( \varphi, \psi) } {( (R +r \cos  \psi) \cos \varphi, (R+ r \cos  \psi ) \sin \varphi ,  r \sin  \psi )
} {}
eine
\definitionsverweis {Bijektion}{}{}
ist.

}
{} {}

\inputaufgabe
{6}
{

Es sei $T$ ein
\definitionsverweis {Torus}{}{}
und seien
\mathl{P,Q \in T}{} zwei Punkte. Zeige, dass es eine gemeinsame Kartenumgebung
\mathl{P,Q \in U \subseteq T}{} derart gibt, dass die Kartenabbildung
\maabbdisp {\alpha} { U } { V
} {}
eine
\definitionsverweis {Homöomorphie}{}{}
mit
\mathl{V= {]0,1[} \times  {]0,1[}}{} ergibt.

}
{} {}

\inputaufgabe
{3}
{

Zeige, dass das
\definitionsverweis {relative Produkt}{}{}
$\R \times_\R \R$ für stetig differenzierbare Funktionen
\maabbdisp {f,g} {\R} {\R
} {}
keine
\definitionsverweis {topologische Mannigfaltigkeit}{}{}
sein muss.

}
{} {}

\inputaufgabe
{2}
{

Es sei

\mavergleichskettedisp
{\vergleichskette
{ D
}
{ =} { \left(  xy^2-  \sin z  \right)  \partial_x +  \left(  e^{x+y+z}  \right) \partial_y+  \left(  xy^3z  \right) \partial_z
}
{ } {
}
{ } {
}
{ } {
}
}
{}{}{.}
Berechne

\mathdisp {D  \left(  7e^{xy}-  \cos \left( z^3+y^5 \right) -x^4y^3 + y^2z^3e^{z}  \right)} { . }

}
{} {}

\inputaufgabe
{2}
{

Es sei

\mavergleichskette
{\vergleichskette
{ W
}
{ \subseteq }{ \R^n
}
{  }{
}
{    }{
}
{   }{
}
}
{}{}{}
\definitionsverweis {offen}{}{,}
\maabb {h} { W } { \R
} {}
eine
\definitionsverweis {stetig differenzierbare Funktion}{}{}
und

\mavergleichskette
{\vergleichskette
{ Y
}
{ = }{ h^{-1}(0)
}
{  }{
}
{    }{
}
{   }{
}
}
{}{}{}
die
\definitionsverweis {Faser}{}{}
zu $0$, wobei $h$ in jedem Punkt von $Y$
\definitionsverweis {regulär}{}{}
sei und somit eine diffferenzierbare Mannigfaltigkeit definiere. Es sei $D$ ein Differentialoperator erster Ordnung zu einem stetigen Vektorfeld auf $Y$. Zeige

\mavergleichskettedisp
{\vergleichskette
{ D(gh)
}
{ =} { 0
}
{ } {
}
{ } {
}
{ } {
}
}
{}{}{}
für jede differenzierbare Funktion $g$ auf $Y$.

}
{} {}

\inputaufgabe
{4 (1+1+1+1)}
{

Wir betrachten auf dem $\R^3$ die
\definitionsverweis {Vektorfelder}{}{}

\mavergleichskettedisp
{\vergleichskette
{ D
}
{ =} { -y \partial_x  + z\partial_y -x \partial_z
}
{ } {
}
{ } {
}
{ } {
}
}
{}{}{,}

\mavergleichskettedisp
{\vergleichskette
{ E
}
{ =} { 5 \partial_x - 3\partial_y +  \cos \left(  xy \right)  \partial_z
}
{ } {
}
{ } {
}
{ } {
}
}
{}{}{,}
und

\mavergleichskettedisp
{\vergleichskette
{ F
}
{ =} { 4e^{xyz} \partial_y -z^3 \partial_z
}
{ } {
}
{ } {
}
{ } {
}
}
{}{}{.}
\aufzaehlungvier{Berechne
\mathl{[D,E]}{.}
}{Berechne
\mathl{[D,F]}{.}
}{Berechne
\mathl{[E,F]}{.}
}{Berechne
\mathl{[ [D,E], F]}{.}
}

}
{} {}

 [[Kategorie:Latexseite]]
